\section{Structural and Cross-Level Properties}
\label{sec:structural-overview}
Appendix~\ref{sec:graph-definitions} defines the dependency graphs at three granularity levels; the detailed per-graph network statistics appear in Appendices~\ref{app:module-detail}--\ref{app:ns-detail}. We now ask two complementary questions: (i)~do the three graphs share common structural signatures, or does each level have its own character? and (ii)~how do the three organizational systems---file boundaries, namespace boundaries, and logical dependencies---align and diverge when overlaid?

The first question is addressed by five network-analytic tools applied uniformly across all three levels: degree distribution (\S\ref{sec:sa-degree}), DAG depth (\S\ref{sec:dag-depth}), centrality (\S\ref{sec:sa-centrality}), community detection (\S\ref{sec:sa-community}), and robustness (\S\ref{sec:sa-robustness}). The full analysis is reported in Appendix~\ref{sec:structural-analysis}; we summarize the key findings here. The second question is addressed by three cross-level measurements: namespace--module alignment, module cohesion, and namespace--declaration interaction. The full analysis is reported in Appendix~\ref{sec:cross_level}; this section extracts the definitions and results that are essential to the main narrative.

\subsection{Structural Properties}
\label{sec:structural-summary}
\label{sec:graph-analysis}

\subsubsection{Per-Graph Findings}

\paragraph{The module graph.}
\label{sec:module_conclusion}
\module{Mathlib}'s module graph is a \textit{deep, engineered backbone} supporting a vast, shallow application layer. Three features stand out:

\begin{enumerate}
    \item \textit{Development-ergonomic redundancy as a feature.} A significant fraction of import edges are transitively redundant (even after \texttt{shake} linting), representing a development ergonomics layer where developers sacrifice logical minimality for usability.

    \item \textit{Funnel topology and compilation bottlenecks.} The layer-width profile reveals a massive leaf layer resting on a deep, narrow foundation, creating recompilation cascades when the narrow ``waist'' is modified---a scalability challenge addressed by Lean's module system~\cite{baanen2025} and distributed builds~\cite{huch2024}.

    \item \textit{Multidimensional centrality.} In-degree, PageRank, and betweenness identify largely disjoint sets of important modules, suggesting that structural importance is irreducibly multidimensional (\S\ref{sec:sa-centrality}).
\end{enumerate}

\noindent
These observations suggest that \module{Mathlib} has evolved into a \textit{curated small-world network} in which the module tree~$T$ and the import graph~$G_{\mathrm{module}}$ continuously reshape each other. A detailed discussion is provided in Appendix~\ref{app:module-analysis}.

\paragraph{The declaration graph.}
\label{sec:thm-conclusion}
The declaration graph dissolves the module boundary, exposing individual logical dependencies. Three findings stand out:

\begin{enumerate}
    \item \textit{Double infrastructure.} The hub structure decomposes into a \emph{language infrastructure} layer determined by Lean's type system and a \emph{mathematical infrastructure} layer determined by mathematical logic (Remark~\ref{rem:two-layers}). A different proof assistant would produce a different hub ranking.

    \item \textit{Emergent disciplinary structure.} Louvain community detection recovers recognizable mathematical disciplines from dependency topology alone, with only moderate alignment to the namespace hierarchy (\S\ref{sec:sa-community}).

    \item \textit{Deep cross-disciplinary entanglement.} Cross-namespace edge density at the declaration level far exceeds the module-level figure. Module boundaries conceal the true extent of mathematical interdependence.
\end{enumerate}

\noindent
Even at this finer resolution, the process leaves indelible traces in the product. A detailed discussion is provided in Appendix~\ref{app:decl-analysis}.

\paragraph{The namespace graph.}
\label{sec:ns-analysis}
The namespace graph is the only level admitting a continuous parameterization (Definition~\ref{def:ns-graph}), yielding a family $G_{\mathrm{ns}}^{(k)}$ indexed by truncation depth. Three findings stand out:

\begin{enumerate}
    \item \textit{Rapid containment decay.} The organizing power of human-imposed boundaries diminishes steeply with granularity, with the sharpest drop at the discipline--sub-discipline transition.

    \item \textit{Structural interpolation.} $G_{\mathrm{ns}}^{(2)}$ consistently interpolates between the module and declaration levels across all structural metrics (\S\ref{sec:sa-community}).

    \item \textit{Infrastructure bottleneck.} The namespace graph is more vulnerable to targeted attack than the declaration graph (\S\ref{sec:sa-robustness}), reflecting the concentration of connectivity in a small number of infrastructure namespaces.
\end{enumerate}

\noindent
A detailed discussion is provided in Appendix~\ref{app:ns-analysis}.

\subsubsection{Cross-Level Structural Comparison}

\paragraph{Degree distributions.}
All three graphs share heavy-tailed degree distributions best fit by lognormal rather than pure power-law models, confirming a universal ``few hubs, many peripheral nodes'' pattern across granularities. The key cross-level invariant is directional: in-degree is open-ended (a product property), while out-degree is bounded (a process property). This asymmetry intensifies at finer resolution---at the declaration level the standard deviation ratio exceeds $21{:}1$ (Tables~\ref{tab:degree-stats}, \ref{tab:thm-degree}, \ref{tab:ns-degree-stats}; Appendix~\ref{sec:sa-degree}).

\paragraph{DAG depth.}
The depth ordering across levels reveals a striking inversion: $153$ layers (module), $84$ (declaration), $8$ (namespace after SCC condensation). Human-organized structure is the deepest; logical structure is the shallowest (Appendix~\ref{sec:dag-depth}). Depth is thus a feature of the production process, not of the mathematical product.

\paragraph{Centrality.}
In-degree, PageRank, and betweenness identify largely disjoint sets of important nodes at every level, confirming that structural importance is irreducibly multidimensional. The divergence sharpens at finer granularity: at the declaration level, no node appears in all three top-$10$ lists (Appendix~\ref{sec:sa-centrality}).

\paragraph{Community detection.}
Louvain modularity decreases from $0.64$ (module) to $0.48$ (declaration) to $0.28$ (namespace), reflecting progressively denser cross-boundary connectivity. Community--namespace alignment follows a parallel decline (NMI from $0.42$ to $0.25$). Only \ns{CategoryTheory} resists this trend, achieving $97.4\%$ community concentration at the declaration level---a unique degree of disciplinary self-containment (Appendix~\ref{sec:sa-community}).

\paragraph{Robustness.}
All three graphs show the classic hub-vulnerability signature---resilient to random failure, fragile under targeted attack---but severity varies: at $20\%$ targeted removal, connectivity drops to $58\%$ (module), $46\%$ (declaration), $21\%$ (namespace). The namespace graph's heightened fragility reflects the concentration of connectivity in a few infrastructure namespaces (Appendix~\ref{sec:sa-robustness}).

\subsection{Cross-Level Alignment}
\label{sec:cross-level-summary}

\subsubsection{Module--Declaration Interaction}
\label{sec:cohesion-main}

In $G_{\mathrm{thm}}$, each module becomes a \emph{region}---a set of declaration-vertices sharing a common source file. For each module~$m$, we partition the incident edges into internal edges $E_{\mathrm{int}}(m)$ (both endpoints in~$m$) and external edges $E_{\mathrm{ext}}(m)$ (exactly one endpoint in~$m$). Figure~\ref{fig:module-region-app} illustrates this partition.

\begin{definition}[Module cohesion]
\label{def:cohesion}
The \emph{cohesion} of a module~$m$ is
$\mathrm{coh}(m) = |E_{\mathrm{int}}(m)| / (|E_{\mathrm{int}}(m)| + |E_{\mathrm{ext}}(m)|)$,
with $\mathrm{coh}(m) = 0$ when both sets are empty. See Appendix~\ref{sec:cohesion} for the full formal development.
\end{definition}

Module cohesion is low ($\bar{h} = 0.107$, median $0.074$), indicating that file boundaries and dependency structure are only weakly aligned. Aggregating $G_{\mathrm{thm}}$ to the module level amplifies the edge count by $9.1\times$ ($215{,}211$ cross-file declaration edges vs.\ $23{,}235$ import edges), confirming that the module graph acts as a sparse gateway to a vast transitive dependency space. Full derivation and discussion appear in Appendix~\ref{sec:cohesion}.

The contrast between module-level and namespace-level grouping is visualised in Figure~\ref{fig:ns-vs-mod-cross} (Appendix~\ref{sec:ns-decl-interaction}): the same declarations and edges yield cohesion~$\approx 10.7\%$ when grouped by file but containment~$\approx 22.2\%$ when grouped by namespace.

\subsubsection{Namespace--Declaration Interaction and Granularity Equilibrium}
\label{sec:ns-decl-main}
\label{sec:granularity-main}

A consistent ordering emerges across boundary types: namespace containment at depth~$1$ ($22.2\%$) exceeds file-level containment ($15.6\%$), which exceeds module cohesion ($10.7\%$). Regardless of which human boundary one draws, the majority of mathematical dependencies cross it. The granularity of modules (${\sim}44$ declarations per module) and namespaces (${\sim}20$ declarations per leaf namespace) is itself a process-trace: a compromise shaped by human cognitive constraints rather than mathematical necessity. The full analysis, including depth asymmetry and the granularity spectrum, is provided in Appendix~\ref{sec:ns-decl-interaction} and~\ref{sec:three-layers}.

\subsubsection{Three Layers Combined}
\label{sec:three-layers-main}

Table~\ref{tab:three-layers-app} (Appendix~\ref{sec:three-layers}) assembles the cross-level measurements. Human organizational decisions---naming and filing---agree with each other (NMI~$= 0.71$) far more than either agrees with the logical structure of the dependency graph (NMI~$= 0.34$, cohesion~$0.107$, containment~$22.2\%$). The divergence is systematic, not accidental: it arises from the fundamental mismatch between the hierarchical nature of human organizational process and the cross-cutting nature of inherited mathematical structure.
